
\documentclass[11pt]{article}
\usepackage{fontspec}
\setmainfont{Times New Roman}
\setsansfont{Arial}
\setmonofont{Consolas}
\usepackage{amsmath,amssymb,mathtools}
\usepackage{geometry}
\usepackage{microtype}
\geometry{a4paper,margin=2.05cm}
\setlength{\parindent}{1.1em}
\setlength{\parskip}{0pt}
\providecommand{\Om}{\Omega}
\providecommand{\N}{\mathrm{N}}
\newcommand{\foreign}[1]{#1}
\emergencystretch=220em
\allowdisplaybreaks
\sloppy
\hbadness=10000
\vbadness=10000
\hfuzz=5pt
\vfuzz=12pt
\raggedbottom
\begin{document}

\section*{38. Zajedno s R. Brauerom i H. Hasse: dokaz glavnoj teoremy v teoriji algebr}

\begin{center}
\foreign{Journal f. d. reine u. angew. Math. 167 (1932), s. 399--404}
\end{center}

\begin{center}
\textbf{Dokaz glavnoj teoremy v teoriji algebr.}\\[0.5em]
R. Brauer v Königsbergu, H. Hasse v Marburgu i E. Noether v Göttingenu\footnote{Sostavjenje toj zapisky vzal na sebe H. Hasse.}.
\end{center}

Nakonec našim zajednym usiljem udalo se dokazati pravdivost slědujuče teoremy, ktora ima osnovno značenje dlja strukturnoj teorije algebr nad algebraičnymi čislovymi poljami, a takože poza toj oblasti:

\paragraph{Glavna teorema.}
\emph{Vsaka normalna dělitvena algebra nad algebraičnym čislovym poljem jest ciklična (abo, kako se takože govori, Dicksonova tipa).}

Za nas jest osobna radost predstaviti toj rezultat, kako uspěh v bitnosti dolžny \(p\)-adičnomu metodu, gospodinu Kurtu Henselu, osnovatelju togo metoda, k jego sedmdesetomu rođenju.

Naš dokaz sostoi iz trěh redukcij; každy iz nas pridal jednu.\footnote{One se daju v porjadku svojego nastanka, protivnom sistematičnomu porjadku.}

\paragraph{1.}
Prvu redukciju dal H. Hasse na osnově nedavno razvitoj jim teorije cikličnyh algebr nad algebraičnymi čislovymi poljami.\footnote{\foreign{H. Hasse, Theorie der zyklischen Algebren über einem algebraischen Zahlkörper, Gött. Nachr. 1931.} Podrobno predstavjenje dokazov, v ktrom se osobno razviva i teorija polj razloženja i skrěženih produktov, razvijena E. Noether v predavanju i tu, kako tam, osnovna, skoro izide v \foreign{Trans. Amer. Math. Soc.} pod naslovom \foreign{Theory of cyclic algebras over an algebraic number field}. Dalje se ta rabota cituje kako H. H, 1--6 sostavljaju, izključno razliky jezika, prvu zapisku.}

\paragraph{Redukcija 1.}
Glavna teorema jest dokazana, ako jest pokazano:
\[
\text{I. Vsaka povsudu razložena algebra nad }\Omega\text{ jest } \sim \Omega .
\]
Radi kratkosti izraz ``algebra \(A\) nad \(\Omega\)'' tu i dalje vsegda znači ``normalna prosta algebra \(A\) nad algebraičnym čislovym poljem \(\Omega\)'', to jest prosty hiperkompleksny sistem s centrom \(\Omega\). Dalej, \(A\sim\Omega\) znači, že \(A\) jest polna matrična algebra nad \(\Omega\), to jest prinadleži specialnoj klasi podobnosti, opreděljenoj samym \(\Omega\) kako dělitvenoj algebra nad \(\Omega\). Nakonec, pod povsudu razloženoju algebra nad \(\Omega\) rozuměje se taka, pri ktoroj za kažno primno město \(\mathfrak p\) polja \(\Omega\) \( \mathfrak p\)-adično razširjenje, to jest algebra nastajuča razširjenjem koeficientnogo polja \(\Omega\) do prinadležnogo \( \mathfrak p\)-adičnogo polja \(\Omega_{\mathfrak p}\), razlaga se:
\[
A_{\mathfrak p}\sim \Omega_{\mathfrak p}.
\]

\paragraph{Dokaz.}
Nehaj \(D\) bude dělitvena algebra nad \(\Omega\). Ako \(Z\) jest ciklično polje nad \(\Omega\) tako, že za kažno primno město \(\mathfrak p\) polja \(\Omega\) \( \mathfrak p\)-stepenj polja \(Z\) jest mnogokratnik \( \mathfrak p\)-indeksa \(D\), togda za prinadležne primne dělitelji \(\mathfrak P\) v \(Z\) kažno \( \mathfrak P\)-adično polje \(Z_{\mathfrak P}\) jest polje razloženja za \(D_{\mathfrak p}\) [H, 18.1]. Iz togo slěduje, že algebra \(D_Z\), nastala razširjenjem koeficientnogo polja \(D\) do \(Z\), povsudu razlaga se nad \(Z\). Ako I uže stoji za vsako \(\Omega\), togda \(D_Z\sim Z\). To govori, že \(D\) ima ciklično polje razloženja \(Z\), teda može byti ciklično predstavljena [H, 5]. Potom \(D\) ima i tako ciklično polje razloženja, čij stepenj sovpada so stepenjem same \(D\); to jest, \(D\) jest ciklična [H, 6, teorema 6]. Pod uslovom I glavna teorema jest tedy dokazana.

\paragraph{2.}
Za drugu redukciju R. Brauer dal odločujuči potisk, kogda v pismu soobčil H. Hasse, kako srodno pytanje o točnoj cěnnosti eksponenta dělitvenoj algebry (ktore vpročem glavna teorema takože rěša, vidi niže) možno s pomočju Sylowovoj grupnoj teoremy svesti na slučaj rěšimogo polja razloženja. Ta mysl potom dala se upotrěbiti i za odgovarajuču dalju redukciju pytanja cikličnosti.

\paragraph{Redukcija 2.}
Tvrdženje I jest dokazano, ako jest pokazano:
\[
\text{II. Vsaka povsudu razložena algebra s rěšimym poljem razloženja nad }\Omega\text{ jest }\sim\Omega .
\]

\paragraph{Dokaz.}
Nehaj \(A\) bude povsudu razložena algebra nad \(\Omega\). Nehaj \(K\) bude Galoisovo polje razloženja za \(A\), \(p\) proizvolno primno čislo, a \(\Sigma\) jedno iz Sylowovyh polj polja \(K\) nad \(\Omega\) odnosimyh do \(p\) (to jest invariantno polje Sylowovoj \(p\)-podgrupy Galoisovoj grupy). Togda \(A_{\Sigma}\) jest takože povsudu razložena algebra nad \(\Sigma\), ktora ima rěšimo polje razloženja \(K\). Ako II uže stoji za vsako \(\Omega\), togda \(A_{\Sigma}\sim \Sigma\). To govori, že \(A\) ima polje razloženja \(\Sigma\) stepenja primnogo k \(p\). Indeks \(A\), kako dělitelj togo stepenja, jest tedy primny k \(p\). Poněže to drži za vsako primno čislo \(p\), indeks jest 1, to jest \(A\sim\Omega\). Pod uslovom II jest tedy I dokazano.\footnote{Mysl redukcije na rěšimo polje razloženja s pomočju Sylowovoj grupnoj teoremy R. Brauer upotrěbil uže raněje, imenno da pokaže, že každy primny dělitelj indeksa pojavja se i v eksponentu (\foreign{Über den Zusammenhang von arithmetischen und invariantentheoretischen Eigenschaften von Gruppen linearer Substitutionen, Berl. Akad.-Ber. 1926}). Nedavno A. A. Albert za tu mysl, i obće za rjad obćih teorem teorije R. Brauera i E. Noether, razvil prosty dokazy nezavisne od teorije predstavjenj: \foreign{1. On direct products, cyclic division algebras, and pure Riemann matrices; 2. On direct products;} oboje v \foreign{Trans. Amer. Math. Soc. 33 (1931)}; za tu redukciju vidi osobno teoremu 23 v 2.

\emph{Dodatok v času pečati.} Dalje A. A. Albert, imějuči pismeno soobčenje H. Hasse, že glavna teorema u njega jest dokazana za abelove algebry (vidi slědujuče v tekstu), iz togo neposrědnje, nezavisno od nas, vyvel slědujuče fakty:

a) glavnu teoremu za stepenje formy \(2^e\),

b) niže slědujuču teoremu 1 (eksponent = indeks),

c) pokraj osnovnoj mysli redukcije 2 takože mysl slědujučej redukcije 3, prirodno bez odnosa k redukciji 1, i sootvětně s rezultatom: za dělitvene algebry \(D\) stepenja \(p^e\) nad \(\Omega\) jest razširjeno polje \(\Omega'\) stepenja primnogo k \(p\) nad \(\Omega\), tako že \(D_{\Omega'}\) jest ciklična.

Vsi tri rezultaty sut, razuměje se, prejděny našim medžutym izvedenym dokazom glavnoj teoremy. Oni pak pokazuju, že i A. A. Albert ima nezavisnu čest v dokazu glavnoj teoremy. Nakonec A. A. Albert (poslě znanja našego dokaza glavnoj teoremy) jošče primětil, že naša centralna teorema I slěduje v nekoliko linijah iz teorem 13, 10, 9 jego raboty v pečati (\foreign{Bull. Amer. Math. Soc. 37 (1931)}). Dokaz tih teorem v bitnosti opira se na te same zaključki kako naše redukcije 2 i 3.}

\paragraph{3.}
Tretju redukciju, ktora potom, kako H. Hasse razpoznal imějuči prvy dvě redukcije, vodi k konečnomu dokazu, dala E. Noether po povodu soobčenja H. Hasse; v njem glavna teorema byla dokazana za specialny slučaj abelovogo polja razloženja, imenno s pomočju obćej redukcije 1 i formulnogo snižanja prinadležnoj faktornoj sistemy. Kako jadro toga E. Noether vydělila právě tretju redukciju. Nezavisno od E. Noether, vpročem, i R. Brauer byl uže promyslil tu tretju redukciju.

\paragraph{Redukcija 3.}
Tvrdženje II jest dokazano, ako jest pokazano:
\[
\begin{gathered}\text{III. Vsaka povsudu razložena algebra s cikličnym poljem razloženja}\\\text{primnogo stepenja nad }\Omega\text{ jest }\sim\Omega .\end{gathered}
\]

\paragraph{Dokaz.}
Nehaj \(A\) bude povsudu razložena algebra s rěšimym poljem razloženja \(K\) nad \(\Omega\). Nehaj
\[
K=\Lambda_0>\Lambda_1>\cdots>\Lambda_r=\Omega
\]
bude veriga polj medžu \(K\) i \(\Omega\) taka, že vsegda \(\Lambda_i\) nad \(\Lambda_{i+1}\) jest ciklično primnogo stepenja. Togda najprvo \(A_{\Lambda_1}\) jest takože povsudu razložena algebra nad \(\Lambda_1\), ktora ima ciklično polje razloženja primnogo stepenja \(K=\Lambda_0\). Ako III uže stoji za vsako \(\Omega\), togda \(A_{\Lambda_1}\sim \Lambda_1\). To govori, že \(\Lambda_1\) jest polje razloženja za \(A\). Potom, počinajuči od \(\Lambda_1\) vměsto \(\Lambda_0\), točno tako dokazuje se \(A_{\Lambda_2}\sim\Lambda_2\), i tako dalje, dokolě nakonec \(A_{\Lambda_r}\sim\Lambda_r\), to jest \(A\sim\Omega\). Pod uslovom III jest tedy II dokazano.

\paragraph{4.}
No III stoji po rezultatam H. Hasse [H, 24]. Tedy nazad čerez redukcije 3, 2, 1 slěduje pravdivost glavnoj teoremy.

E. Noether pri tom jošče ukazuje slědujuče.
Pravdivost III, obćeje za ciklična polja razloženja proizvolnogo stepenja, svodi se na normnu teoremu Hasse\footnote{\label{fn:p38-hasse-normensatz}Navedeno v H, 3.11; dokazano v: \foreign{H. Hasse, Beweis eines Satzes und Widerlegung einer Vermutung über das allgemeine Normenrestsymbol, Gött. Nachr. 1931}.}. Za redukciju 3 potrěbny jest ale samo specialny slučaj primnogo stepenja, to jest izvorna normna teorema Hilbert--Furtwängler\footnote{Vidi \foreign{H. Hasse, Bericht über neuere Untersuchungen und Probleme in der Theorie der algebraischen Zahlkörper II, Jahresber. der D. M.-V., Erg.-Bd. 6 (1930), § 8}, i tam navedenu literaturu.}. Tak redukcija III daje novy prosty dokaz normnoj teoremy Hasse. Taj dokaz šel by tako:

Nehaj \(Z\) bude ciklično polje nad \(\Omega\), a \(\alpha\) čislo iz \(\Omega\), za ktore normno-ostatkovy symbol
\[
\left(\frac{\alpha,Z}{\mathfrak p}\right)=1
\]
jest za kažno primno město \(\mathfrak p\) polja \(\Omega\). Togda vsaka ciklična algebra
\[
A=(\alpha,Z)
\]
\footnote{V označenju H, 1. Podanje avtomorfizma \(S\) polja \(Z\), svjazanogo s \(\alpha\), bylo tu opuščeno kako nesuštno.} povsudu razlaga se [H, 17.7]. Po redukciji 3 tedy, upotrěbivši samo normnu teoremu Hilbert--Furtwängler, slěduje \(A\sim\Omega\). Potom ale \(\alpha\) jest norma elementa iz \(Z\) [H, 15.4].

V svjazi s normnoju teoremoju dalje priměčamo, že teoremu I treba smatrati kako jej pravilno obobčenje na vyšje, takože neabelove, slučaje, medžutym kako doslovno obobčenje, kako H. Hasse pokazal\textsuperscript{\ref{fn:p38-hasse-normensatz}}, ne jest obće pravilno.

\paragraph{5.}
Iz glavnoj teoremy odmah slěduje odgovor na pytanje, k kotoromu izvorno odnosila se druga redukcijna mysl:

\paragraph{Teorema 1.}
Eksponent normalnoj prostoj algebry nad algebraičnym čislovym poljem ravny jest jej indeksu.

Dalje iz teoremy I slěduje znatny fakt:

\paragraph{Teorema 2.}
Osnovny ideal pravej normalnoj dělitvenoj algebry nad \(\Omega\) dělimo jest najmenej jednym primnym městom \(\Omega\), i čak najmenej dvěma.

Pri tom osnovny ideal opredělja se napr. kako reducirana norma reduciranoj diferenty.\footnote{V specialnom slučaju racionalnyh kvaternionnyh algebr to svodi se na pojam ``osnovnogo čisla'', vvedeny H. Brandtom (\foreign{Idealtheorie in Quaternionenalgebren, Math. Ann. 99 (1928)}). Poněže tu ide ne o čislo, ale o ideal, trebalo je govoriti ``osnovny ideal''; to ne smije se měšati s Dedekindovym označenjem osnovny ideal i osnovno čislo za diferentu i diskriminantu, ktore právě iz navedene pričiny za relativna polja pokazyva se neupotrěbno. O definiciji reduciranoj diferenty vidi slědujuču notu. Reduciranu normu ideala E. Noether takože opredělja ``po primnyh městah'', imenno, sootvětně faktu, že za pojedina primna města každy ideal jest glavni ideal, prosto tvorjenjem reducirovanyh čislovyh norm čisel bazisa glavnogo ideala za pojedina primna města.} K tomu beskonečno primno město vključa se v osnovny ideal togda i samo togda, kogda algebra za njega reduciruje se na kvaternionnu algebru, a ne na realno abo kompleksno čislovo polje.

\paragraph{Dokaz.}
Po rezultatam H. Hasse\footnote{\label{fn:p38-hasse-padisch}\foreign{H. Hasse, Über \(p\)-adische Schiefkörper und ihre Bedeutung für die Arithmetik hyperkomplexer Zahlsysteme, Math. Ann. 104 (1931)}; vidi tam teoremy 42, 59.} primno město \(\mathfrak p\) polja \(\Omega\) vhodi v osnovny ideal togda i samo togda, kogda \( \mathfrak p\)-indeks različny jest od \(1\). Ako by vsi \( \mathfrak p\)-indeksy byli ravne \(1\), algebra razlagala by se povsudu, i po teoremě I ne byla by prava dělitvena algebra. (Že ni jeden samy \( \mathfrak p\)-indeks ne može byti različny od \(1\), slěduje iz togo, že za normno-ostatkove symbole, čiji porjadki sut \( \mathfrak p\)-indeksy [H, 17.7], drži produktna teorema, to jest zakon vzajemnosti [H, 3.8].)

\paragraph{6.}
Dalje teorema I umožnja opreděliti, aritmetično označiti, vsaka polja razloženja \(K\), prinadležna algebri \(A\) nad \(\Omega\). V točnom obobčenju stanja, ktore H. Hasse uže ustanovil v specialnom slučaju cikličnyh \(A,K\) [H, 6, teorema 2], dobivamo:

\paragraph{Teorema 3.}
Za algebru \(A\) nad \(\Omega\) algebraično čislovo polje \(K\) nad \(\Omega\) jest polje razloženja togda i samo togda, kogda za vse primne dělitelji \(\mathfrak P_i\) v \(K\) vseh primnyh měst \(\mathfrak p\) polja \(\Omega\) sootvětny \( \mathfrak P_i\)-stepenj \(n_{\mathfrak P_i}\) polja \(K\) jest mnogokratnik \( \mathfrak p\)-indeksa \(m_{\mathfrak p}\) algebry \(A\).

Pri tom \( \mathfrak P_i\)-stepenj polja \(K\) jest stepenj prinadležnogo \( \mathfrak P_i\)-adičnogo polja \(K_{\mathfrak P_i}\) nad \( \mathfrak p\)-adičnym poljem \(\Omega_{\mathfrak p}\), to jest, ako
\[
\mathfrak p=\prod_i \mathfrak P_i^{\,e_{\mathfrak P_i}},
        \qquad N_{K/\Omega}(\mathfrak P_i)=\mathfrak p^{\,f_{\mathfrak P_i}},
\]
jest razloženje \(\mathfrak p\) v \(K\), produkt stepenja \(f_{\mathfrak P_i}\) i porjadka razvětvjenja \(e_{\mathfrak P_i}\) dělitelja \(\mathfrak P_i\) glede \(\mathfrak p\), \(n_{\mathfrak P_i}=f_{\mathfrak P_i}e_{\mathfrak P_i}\) [H, 4]. A \( \mathfrak p\)-indeks \(A\) jest indeks \(m_{\mathfrak p}\) \( \mathfrak p\)-adičnogo razširjenja \(A_{\mathfrak p}\) [H, 6].

\paragraph{Dokaz.}
Že \(K\) jest polje razloženja za \(A\), to jest \(A_K\sim K\), po teoremě I ravnoznačno jest tomu, že \(A_{K_{\mathfrak P_i}}\sim K_{\mathfrak P_i}\) za každy \(\mathfrak P_i\).

Za konečno primno město \(\mathfrak p\) nehaj
\[
        A_{\mathfrak p}\sim(\pi,W_{\mathfrak p})
\]
bude aritmetično odlično ciklično predstavjenje \(A_{\mathfrak p}\) [H, 16; k tomu l. c.\textsuperscript{\ref{fn:p38-hasse-padisch}}, teorema 38], tedy \(W_{\mathfrak p}\) nerazvětvjeno polje stepenja \(m_{\mathfrak p}\) nad \(\Omega_{\mathfrak p}\), a \(\pi\) čislo iz \(\Omega_{\mathfrak p}\), točno dělimo s \(\mathfrak p^1\). Dalej nehaj \(K_{\mathfrak P_i}W_{\mathfrak p}\) bude kompozitum i
\[
        \Delta_{\mathfrak P_i}=K_{\mathfrak P_i}\cap W_{\mathfrak p}
\]
prerěz \(W_{\mathfrak p}\) i \(K_{\mathfrak P_i}\). Poněže največje nerazvětvjeno podpolje v \(K_{\mathfrak P_i}\) ima stepenj \(f_{\mathfrak P_i}\) nad \(\Omega_{\mathfrak p}\), i poněže nad \(\Omega_{\mathfrak p}\) za každy stepenj jest samo jedno nerazvětvjeno polje, \(\Delta_{\mathfrak P_i}\)
\[
        \text{ima stepenj }d_{\mathfrak P_i}=(m_{\mathfrak p},f_{\mathfrak P_i})\text{ nad }\Omega_{\mathfrak p}.
\]
Zato \(K_{\mathfrak P_i}W_{\mathfrak p}\) ima stepenj \(m_{\mathfrak p}/d_{\mathfrak P_i}\) nad \(K_{\mathfrak P_i}\), pri tom nerazvětvjen.

Teper
\[
        A_{K_{\mathfrak P_i}}=(A_{\mathfrak p})_{K_{\mathfrak P_i}}\sim(\pi,K_{\mathfrak P_i}W_{\mathfrak p})
\]
[H, 15.5]. Sledovateljno \(A_{K_{\mathfrak P_i}}\sim K_{\mathfrak P_i}\) ravnoznačno jest tomu, že \(\pi\) jest norma iz \(K_{\mathfrak P_i}W_{\mathfrak p}\) glede \(K_{\mathfrak P_i}\). No zbog nerazvětvjenosti \(K_{\mathfrak P_i}W_{\mathfrak p}\) nad \(K_{\mathfrak P_i}\) to biva togda i samo togda, kogda porjadkovo čislo \(e_{\mathfrak P_i}\) od \(\pi\) v \(\mathfrak P_i\) jest mnogokratnik stepenja \(m_{\mathfrak p}/d_{\mathfrak P_i}\) polja \(K_{\mathfrak P_i}W_{\mathfrak p}\) nad \(K_{\mathfrak P_i}\), abo, što zbog
\[
\left(\frac{m_{\mathfrak p}}{d_{\mathfrak P_i}},\frac{f_{\mathfrak P_i}}{d_{\mathfrak P_i}}\right)=1
\]
jest to samo, kogda
\[
\frac{f_{\mathfrak P_i}}{d_{\mathfrak P_i}}e_{\mathfrak P_i}=\frac{n_{\mathfrak P_i}}{d_{\mathfrak P_i}}
\]
jest mnogokratnik \(m_{\mathfrak p}/d_{\mathfrak P_i}\), to jest kogda \(n_{\mathfrak P_i}\) jest mnogokratnik \(m_{\mathfrak p}\), kako tvrdženo.

Za slučaj, že \(\mathfrak p\) jest beskonečno primno město i ne trivialny slučaj \(m_{\mathfrak p}=1\), imamo \(m_{\mathfrak p}=2\) i \(A_{\mathfrak p}\) jest kvaternionna algebra nad realnym čislovym poljem \(\Omega_{\mathfrak p}\). Tvrdženje \(A_{K_{\mathfrak P_i}}=(A_{\mathfrak p})_{K_{\mathfrak P_i}}\sim K_{\mathfrak P_i}\) togda ravnoznačno jest tomu, že \(K_{\mathfrak P_i}\) jest kompleksno čislovo polje, to jest \(n_{\mathfrak P_i}=f_{\mathfrak P_i}=2=m_{\mathfrak p}\) (\(e_{\mathfrak P_i}\) tu ne vhodi), što znovu daje tvrdženje (\(n_{\mathfrak P_i}\), kako \(m_{\mathfrak p}\), može iměti samo vrědnosti 1 abo 2).

\paragraph{7.}
Do značečih rezultatov dohodimo, ako pytanje v teoremě 3 o vsih poljah razloženja \(K\) fiksnoj algebry \(A\) obratimo: pri fiksnom algebraičnom polju \(K\) nad \(\Omega\) smatrimo vsotu algebr \(A\) nad \(\Omega\), ktore \(K\) razlaga. Te algebry \(A\) tvore podgrupu \(\mathfrak K\), opreděljenu \(K\), R. Brauerovoj grupy \(\mathfrak A\) vseh algebr, točnije vseh klas podobnyh algebr, nad \(\Omega\).\footnote{\foreign{R. Brauer, Über Systeme hyperkomplexer Größen, Jahresber. d. D. M.-V. 38 (1929), s. 47/48.} Vidi takože H, 13.1.} Ako predpostavimo, že \(K\) jest Galoisovo nad \(\Omega\), dohodimo do teorem, ktore treba smatrati kako obobčenje glavnyh teorem teorije klasnyh polj, teorije relativno-abelovyh čislovyh polj, na obća relativno-Galoisova čislova polja:

\paragraph{Teorema 4.}
\emph{(Teorema razloženja).} Relativny stepenj \(f\) primnyh děliteljev v \(K\) primnogo ideala \(\mathfrak p\) polja \(\Omega\), ktory ne vhodi v relativnu diskriminantu \(K\), ravny jest najranějšemu eksponentu, za ktory
\[
        A_{\mathfrak p}^{\,f}\sim\Omega_{\mathfrak p}
\]
stane za vse algebry \(A\) iz grupy \(\mathfrak K\), pridruženej k \(K\).

\paragraph{Dokaz.}
Poněže \(f\) jest \( \mathfrak p\)-stepenj polja \(K\), jednak za vse primne dělitelji \(\mathfrak p\) v \(K\), a s drugoj strany \( \mathfrak p\)-indeks \(A\) ravny jest indeksu, tedy eksponentu, \(A_{\mathfrak p}\), po teoremě 3 \(f\) vsakako jest mnogokratnik onogo najranějšego eksponenta; za dokaz dostaje pokazati, že v grupě \(\mathfrak K\) sut algebry \(A\) s točnym \( \mathfrak p\)-indeksom \(f\).

Po Frobeniusovoj teoremě gustoty sigurno jest jošče jeden primny ideal \(\mathfrak p'\) v \(\Omega\), ne vhodeči v relativnu diskriminantu \(K\), čiji primni dělitelji v \(K\) iměju relativny stepenj \(f\). Nehaj \(Z\) bude ciklično polje nad \(\Omega\), čij \( \mathfrak p\)-stepenj i \( \mathfrak p'\)-stepenj imata vrědnost \(f\). Dalej nehaj \(\alpha\) bude čislo v \(\Omega\), za ktore normno-ostatkove symbole
\[
\left(\frac{\alpha,Z}{\mathfrak p}\right)\quad\text{i}\quad
\left(\frac{\alpha,Z}{\mathfrak p'}\right)
\]
imaju recipročne vrědnosti porjadka \(f\), dokolě za vse druge dělitelji \(\mathfrak q\) provodnika \(Z\) drži
\[
\left(\frac{\alpha,Z}{\mathfrak q}\right)=1.
\]
Po obobčenoj teoremě o aritmetičnoj progresiji možno k tomu izbrati \(\alpha\) tako, že kromě možebno \(\mathfrak p,\mathfrak p'\) i \(\mathfrak q\) ono ima samo ešte jeden primny ideal \(\mathfrak r\) polja \(\Omega\), točno v prvoj potenciji. Za tego potom po produktnoj teoremě za normno-ostatkovy symbol, zakonu vzajemnosti, takože
\[
\left(\frac{\alpha,Z}{\mathfrak r}\right)=1
\]
[H, 3.8--10]. Vsaka algebra
\[
        (\alpha,Z)=A
\]
ima togda \( \mathfrak p\)-indeks i \( \mathfrak p'\)-indeks \(f\), a za vse druge primne města \(\Omega\) indeks \(1\) [H, 17.7]. Po teoremě 3, s ogledom na predpostavku o \(\mathfrak p\) i izbor \(\mathfrak p'\), slěduje, že \(K\) jest polje razloženja za \(A\). Tak sut vpravdu v grupě \(\mathfrak K\) pokazane algebry \(A\) s \( \mathfrak p\)-indeksom \(f\).

\paragraph{Teorema 5.}
\emph{(Teorema jednoznačnosti i uporjadženja).} Ako \(K\leqq K'\), togda \(\mathfrak K\leqq\mathfrak K'\), i obratno.

Osobno, pridruženje algebrnyh grup \(\mathfrak K\) k Galoisovym poljam \(K\) jest obratno jednoznačno.

\paragraph{Dokaz.}
a) Iz \(K\leqq K'\) trivialno slěduje \(\mathfrak K\leqq\mathfrak K'\); bo vsaka algebra \(A\), razložena od \(K\), a fortiori jest algebra razložena od \(K'\).

b) Obratno nehaj \(\mathfrak K\leqq\mathfrak K'\). Togda po teoremě razloženja za každy nedělitelj \(\mathfrak p\) relativnoj diskriminanty \(f_{\mathfrak p}\mid f'_{\mathfrak p}\). Osobno \(f_{\mathfrak p}=1\) za skoro vse \(\mathfrak p\) s \(f'_{\mathfrak p}=1\). Iz togo po običnom analitičnom zaključnom postupku\footnote{Vidi k tomu napr. H. Hasse [l. c. nota 6], § 25, III.} slěduje \(K\leqq K'\).

\paragraph{8.}
Nakonec izložimo jošče, že glavna teorema daje bitny postup za pytanje, obrabotano I. Schurom\footnote{\foreign{I. Schur, Arithmetische Untersuchungen über endliche Gruppen linearer Substitutionen, Berl. Akad.-Ber. 1906}.}, o čislovyh poljah, v ktoryh možne sut absolutno-nerazložime predstavjenja konečnoj grupy:

\paragraph{Teorema 6.}
Vsa absolutno-nerazložima predstavjenja konečnoj grupy \(\mathfrak G\) možne sut v kružnyh poljah, napr. vsakako v polju \(n^h\)-tyh korenjev jednoty, ako \(n\) jest porjadok \(\mathfrak G\), a \(h\) dostatočno veliko.

\paragraph{Dokaz.}
Prehodjači k racionalno-čislovomu grupovomu kolcu \(G\) grupy \(\mathfrak G\), absolutno-nerazložima predstavjenja \(\Gamma_i\) grupy \(\mathfrak G\) stavaju absolutno-nerazložimymi predstavjenjami prostyh sostavnyh česti \(G_i\) poluprostoj algebry \(G\), i jih centry sut sootvětna polja \(\Omega_i\) harakterov,\footnote{Vidi k tomu: a) \foreign{R. Brauer und E. Noether, Über minimale Zerfällungskörper irreduzibler Darstellungen, Berl. Akad.-Ber. 1927; § 1.} b) \foreign{R. Brauer, Über Systeme hyperkomplexer Zahlen, Math. Zeitschr. 30 (1929);} teorema 3. c) \foreign{E. Noether, Hyperkomplexe Größen und Darstellungstheorie, Math. Zeitschr. 30 (1929); §§ 21, 24, 26.}} zaradi konečnosti \(\mathfrak G\) tedy svakako kružna polja, i sigurno podpolja polja \(n\)-tyh korenjev jednoty.

Po teoremě 3 ciklično polje \(Z_i\) nad \(\Omega_i\) jest polje razloženja za \(G_i\), ako za kažno primno město \(\mathfrak p\) polja \(\Omega_i\) jego \( \mathfrak p\)-stepenj \(n_{i\mathfrak p}\) jest mnogokratnik \( \mathfrak p\)-indeksa \(m_{i\mathfrak p}\) od \(G_i\). \(m_{i\mathfrak p}\) različny jest od 1 samo za primne dělitelji osnovnogo ideala \(G_i\) glede \(\Omega_i\), teda sigurno samo za primne dělitelji absolutnoj diskriminanty \(G_i\); poněže ona vhodi v \(n^n\), a \(n^n\) jest diskriminanta jednogo nemaksimalnogo porjadka v \(G\)\footnote{Vidi E. Noether [l. c. nota 12c], § 26.}, najviše za primne dělitelji \(\mathfrak p\) od \(n\). Da se \(n_{i\mathfrak p}\) za ta \(\mathfrak p\) učini mnogokratnikom \(m_{i\mathfrak p}\), dostaje predpisati, že dotična \(\mathfrak p\) v \(Z_i\) razvětvjena sut porjadkom, každy raz děliteljnym s \(m_{i\mathfrak p}\). To, kako lahko se vidi, daje polje \(n^h\)-tyh korenjev jednoty za dostatočno vysoko \(h\).

V poslědnoj teoremě navedenej raboty I. Schur ustanavlja, že v vsih doselě znanyh slučajah dostaje uže polje \(n\)-tyh korenjev jednoty. Či to vsegda jest tako i či tu razvite metody dostaju za rěšenje togo pytanja, ostaje predmetom dalših izslědovanj.

\begin{center}
Prišlo 11. novembra 1931.
\end{center}

\end{document}
